% !TEX root = ../main.tex
\chapter{Construction of the tangent bundle}\label{ch:l10}
\begin{paracol}{2}

\begin{sources}
\begin{tabularx}{\linewidth}{@{}l l L@{}}
\toprule
\textbf{Video} & \textbf{Length} & \textbf{Content}\\
\midrule
\seg{10}{1}{Lecture 10: Construction of the tangent bundle} & 1:48:50 & Cotangent spaces and tensor spaces at a point; canonical and non-canonical isomorphisms; the gradient $\mathrm d_pf$ as a co-vector and the level sets; the chart-induced dual basis $\mathrm dx^a$;
push-forward (the derivative of a smooth map) and pull-back; vectors are pushed forward, forms pulled back; immersions and embeddings with the circle as example; the Whitney embedding and immersion theorems and the Klein bottle;
the tangent bundle as a set bundle, a correction about a bundle of real and complex lines, the topology and the smooth atlas on $TM$ (compatibility left to the problem sheet).\\
\bottomrule
\end{tabularx}

\smallskip
\textbf{Transcript} \file{materials/transcripts/C10-V01.txt} (YouTube automatic captions).
\end{sources}

\begin{reading}[Where this lecture sits in the other notes]
\emph{Gravity and Light}, Lecture~6 (Fields); \emph{Theoretische Mechanik}, Lecture~5. The tangent bundle and the push-forward are the prerequisites of vector fields; the Whitney theorems, the immersion/embedding distinction and the discussion of the gradient
as a co-vector are, as far as the titles say, specific to this course; the modules announced here are the subject of \cref{ch:l11}.
(From the lecture titles and the plan of the course; not checked in detail.)
\end{reading}

% ==================================================================
\section{Plan of the lecture}\label{sec:tb-plan}

Today some technology is developed, collecting along the way a number of interesting results; the two most important ones \ts{10}{1}{0:04}. First the \textbf{Whitney embedding theorems}, which clarify whether there is really a difference between the intrinsic definition of a manifold and of its tangent spaces and the \emph{extrinsic} one, not
followed in this course (\cref{int:df-embedding}), where the manifold is thought of as embedded in some $\R^N$ and the tangent spaces are planes in the embedding space. The intrinsic point of view was argued for, since for physics, if the manifold is to represent the universe, it would be funny to think of an even bigger space into which the
universe is embedded; the question is whether something could be gained by taking the embedding approach, and the Whitney theorems say that the two are equivalent \ts{10}{1}{0:04} \ts{10}{1}{0:26} \ts{10}{1}{0:47} \ts{10}{1}{1:07}.
Second, the construction of the tangent bundle: vector fields will be defined, but the space of all vector fields is \emph{not} a vector space (the space of all vectors at a point was): it is merely a \emph{module}, a vector space over a commutative ring; modules are the subject of the next lecture (\cref{ch:l11}) \ts{10}{1}{1:28} \ts{10}{1}{1:49} \ts{10}{1}{2:10}.

% ==================================================================
\section{Cotangent spaces and tensor spaces at a point}\label{sec:tb-cotangent}

Having constructed the tangent spaces last time (\cref{ch:l09}) it is quite simple to come up with the cotangent space \ts{10}{1}{2:10} \ts{10}{1}{2:30}.
\begin{definition}[New concepts: cotangent space]\label{def:tb-cotangent}
Let $M$ be a smooth manifold (with its topology and smooth atlas). The \textbf{cotangent space} to $M$ at $p\in M$ is the vector space dual of the tangent space (\cref{def:te-hom}),
\[
  T^*_pM:=(T_pM)^*,
\]
with the inherited addition and scalar multiplication; the star ``jumps'' in the notation. It is a vector space, ``no secret here'' \ts{10}{1}{2:30} \ts{10}{1}{2:51} \ts{10}{1}{3:13} \ts{10}{1}{3:34}.
\end{definition}
\begin{remark}[Finite dimension, and a non-canonical isomorphism]\label{rem:tb-noncanonical}
The manifolds defined in this course are finite dimensional, since one always maps into some $\R^D$; the definition did not include infinite dimensional manifolds, which can be seen as a defect: with more trouble (``not much more, but some more'') one could. Since $\dim T_pM=\dim M<\infty$ (\cref{thm:tv-dim}),
$T_pM$ and $T^*_pM$ have the same dimension and are \emph{isomorphic as vector spaces}: there exists a bijective linear map between them. ``Which one? Many; but none of them can be constructed without extra structure''. The isomorphisms exist but are \textbf{non-canonical}: from the bare vector space structure
one cannot provide the map without picking, say, arbitrarily the entries of a matrix defining it. With extra structure (an inner product, for instance) one can construct one. This is quite different from $V\cong V^{**}$ (\cref{thm:te-dd}, on the current problem sheet): there a \textbf{canonical} isomorphism exists, one that needs no extra structure beyond the level
at which one is discussing, here the level of bare vector spaces \ts{10}{1}{3:55} \ts{10}{1}{4:36} \ts{10}{1}{4:58} \ts{10}{1}{5:38} \ts{10}{1}{5:59} \ts{10}{1}{6:40}.
\end{remark}
For the tangent space there was a nuts-and-bolts construction with curves and directional derivative operators, step by step; for the cotangent space this is not necessary, one just takes the vector space dual. Hence at a point $p$ one has the vector space $T_pM$ (the only thing that had to be built from scratch, over which everything is then pure
linear algebra) and its dual \ts{10}{1}{7:01} \ts{10}{1}{7:21}.
\begin{definition}[New concepts: tensors at a point, $T^r_s(T_pM)$]\label{def:tb-tensors}
Similarly one can construct the space of $(r,s)$ tensors at $p$,
\[
  T^r_s(T_pM):=\{\,T\colon\underbrace{T^*_pM\times\dots\times T^*_pM}_{r}\times\underbrace{T_pM\times\dots\times T_pM}_{s}\to\R\ \text{multilinear}\,\},
\]
in the notation of \cref{def:te-tspace}, with the tangent space $T_pM$ as the underlying vector space ($T$ the tensor, $T_pM$ the tangent space: different $T$'s; no smoothness condition is required, it is a pure vector space matter, the multilinear maps are on vector spaces at the same point). It is a vector space with the pointwise
addition and scalar multiplication, and every tensor at $p$ lives there \ts{10}{1}{7:21} \ts{10}{1}{7:42} \ts{10}{1}{8:23} \ts{10}{1}{8:44} \ts{10}{1}{9:27} \ts{10}{1}{9:47}.
\end{definition}
The only time one has to construct something from scratch is the tangent space itself with the curves; everything else is generated by pure linear algebra. ``Tensors have nothing to do with differential geometry, tensors have all to do with linear algebra and multilinear maps; the only thing we had to construct was an underlying vector space over which the tensor space is
built'', always at this one point $p$ \ts{10}{1}{10:07} \ts{10}{1}{10:29} \ts{10}{1}{10:50}.

% ==================================================================
\section{The gradient}\label{sec:tb-gradient}

\begin{definition}[New concepts: gradient operator, gradient]\label{def:tb-gradient}
Let $f$ be a smooth function on $M$. At every point $p\in M$ there is a linear map, written $\mathrm d_p$ (``if you want to emphasise that we look at this only at one point''; a different map for every point),
\[
  \mathrm d_p\colon C^\infty(M)\to T^*_pM,\qquad f\mapsto\mathrm d_pf,\qquad(\mathrm d_pf)(X):=X(f)\quad\text{for }X\in T_pM.
\]
The map $\mathrm d_p$ is called the \textbf{gradient operator at $p$} and $\mathrm d_pf$ the \textbf{gradient of $f$ at $p$} \ts{10}{1}{10:50} \ts{10}{1}{11:31} \ts{10}{1}{11:52} \ts{10}{1}{12:34} \ts{10}{1}{12:55} \ts{10}{1}{14:42}.
\end{definition}
This seems formalism taken to an extreme: one could simply write $X(f)$ instead of defining a map $\mathrm d_pf\colon T_pM\to\R$ acting on $X$. The reason is that one wants to speak of this object \emph{without providing the vector it acts on}; as in analysis one often wants to talk about the function as such and not about its value $f(x)$ at some point
\ts{10}{1}{13:18} \ts{10}{1}{13:39} \ts{10}{1}{14:00}.
\begin{remark}[The gradient is a co-vector, not a vector]\label{rem:tb-covector}
(a) $\mathrm d_pf\in T^*_pM$ is a \emph{co-vector}, not a vector \ts{10}{1}{14:42} \ts{10}{1}{15:05}. One often hears that the gradient is orthogonal to the level lines of $f$; but orthogonality needs an inner product, and with a bare vector space, which is the case here, there is none in sight.\par
(b) Let $X\in T_pM$ be tangent to the \textbf{level set} $N_c(f):=\{p\in M\mid f(p)=c\}$, $c\in\R$ (``tangent to a level set'' has not been defined formally yet, but the idea: a function on a three-dimensional manifold constant on a sphere, where the tangent to the level set is the tangent to a smooth submanifold; to be treated later). Then
automatically $\mathrm d_pf(X)=0$: the co-vector acts on the vector and annihilates it. The statement ``the gradient is orthogonal to the level sets'' needs no inner product \ts{10}{1}{15:05} \ts{10}{1}{15:25} \ts{10}{1}{16:08} \ts{10}{1}{16:31} \ts{10}{1}{17:12} \ts{10}{1}{17:32}.
In standard vector analysis one commits a ``perversion'': with a metric one constructs a non-canonical isomorphism to raise the index and make the gradient a vector, and then uses the metric again to dot this gradient vector into another vector; the inverse metric used for the first step and the metric used for the second cancel. ``It is totally perverse,
because the gradient naturally is a co-vector'' \ts{10}{1}{17:32} \ts{10}{1}{17:53} \ts{10}{1}{18:14} \ts{10}{1}{18:35}.
\end{remark}

% ==================================================================
\section{The chart-induced dual basis $\mathrm dx^a$}\label{sec:tb-dx}

The gradient is useful in itself, but it can also be used to construct a basis of the cotangent space \ts{10}{1}{19:03} \ts{10}{1}{19:24}.
\begin{definition}[New concepts: chart-induced co-vector basis]\label{def:tb-dx}
Let $p\in M$ and $(U,x)$ a chart with $p\in U$. Each component function $x^a\colon U\to\R$, $a=1,\dots,\dim M$, of the chart map is smooth (\cref{rem:tv-legal}); one can act on it with the gradient operator and obtains the $\dim M$ co-vectors
\[
  \mathrm d_px^a\in T^*_pM,\qquad a=1,\dots,\dim M,
\]
written $\mathrm dx^a|_p$ for short. They are called the \textbf{chart-induced co-vector basis} at $p$ (Theorem below shows that it is a basis, and indeed a dual one) \ts{10}{1}{19:24} \ts{10}{1}{20:05} \ts{10}{1}{20:25} \ts{10}{1}{21:09} \ts{10}{1}{21:29} \ts{10}{1}{22:11}.
\end{definition}
There are enough of them (the cotangent space has dimension $\dim M$ since the manifold is finite dimensional), so they are a candidate for a basis; this is one of those ``definition-theorems'' where one has to provide a proof.
\begin{theorem}[The dual basis]\label{thm:tb-dualbasis}
$\mathrm dx^a|_p(\partial/\partial x^b|_p)=\delta^a_b$. Hence $\mathrm dx^1|_p,\dots,\mathrm dx^{\dim M}|_p$ is a basis of $T^*_pM$, namely the \emph{dual basis} (\cref{def:te-dualbasis}) of the chart-induced basis $\partial/\partial x^b|_p$ of $T_pM$ \ts{10}{1}{22:11} \ts{10}{1}{27:27} \ts{10}{1}{27:48} \ts{10}{1}{28:08}.
\end{theorem}
\begin{proof}
Let the $a$-th co-vector act, as a linear map $T_pM\to\R$, on the $b$-th basis vector. By the definition of $\mathrm d_p$ and then of $\partial/\partial x^b|_p$ (\cref{def:tv-ddx}):
\[
  \mathrm dx^a|_p\Bigl(\frac{\partial}{\partial x^b}\Bigr|_p\Bigr)=\frac{\partial}{\partial x^b}\Bigr|_p(x^a)=\partial_b\bigl(x^a\circ x^{-1}\bigr)\bigl(x(p)\bigr)=\delta^a_b,
\]
since $x^a\circ x^{-1}$ is the projection to the $a$-th component of $\R^{\dim M}$, whose $b$-th partial derivative is $1$ if $b=a$ and $0$ otherwise (``it doesn't matter where you evaluate it''). ``It's also good practice to calculate this: a mere application of the definitions'' \ts{10}{1}{22:11} \ts{10}{1}{23:14} \ts{10}{1}{24:17} \ts{10}{1}{24:58} \ts{10}{1}{26:01} \ts{10}{1}{27:05}.
For linear independence (and the generating property): the equation $\mu_a\,\mathrm dx^a|_p=0$ evaluated on $\partial/\partial x^b|_p$ gives $\mu_b=0$; moreover for $\omega\in T^*_pM$ one has $\omega=\omega(\partial/\partial x^a|_p)\,\mathrm dx^a|_p$ (``for every basis element here one finds one that does not annihilate it and all the others annihilate it''),
and there are $\dim M$ of them: a basis, which by $\delta^a_b$ is the dual one. This proves more than was wanted: it is the dual basis of the dual of $T_pM$, ``I chose this more arcane formulation to use the word dual twice in one sentence, with slightly different meanings'': the dual of $T_pM$ is $T^*_pM$, and the dual basis refers to
applying one basis to the other and getting $\delta^a_b$ \ts{10}{1}{27:27} \ts{10}{1}{27:48} \ts{10}{1}{28:08} \ts{10}{1}{28:29} \ts{10}{1}{28:49}.
\end{proof}
\begin{intuition}[What $\mathrm dx^a$ is: not a small change in $x^a$]\label{int:tb-dx}
It is important not to look at $\mathrm dx^1$ as one does in undergraduate physics, as ``a small change in $x^1$''; it is precisely the co-vector defined above, and that makes full sense; one then has full technical control over what one is doing. Whenever one meets $\mathrm dx^1,\dots,\mathrm dx^{\dim M}$ in the future, for instance in the measure of multiple integrals $\mathrm dx^1\,\mathrm dx^2\,\mathrm dx^3$
(``there is a wedge product between them, we'll come to that''), this is exactly what appears under the integral, and it is not just a symbol: it can be dissected in this way, and it is very important to do so if one wants to do analysis on a smooth manifold \ts{10}{1}{28:49} \ts{10}{1}{29:11} \ts{10}{1}{29:32} \ts{10}{1}{29:53} \ts{10}{1}{30:13}.
\end{intuition}
\begin{exercise}[Components of $\mathrm d_pf$]\label{ex:tb-df}
\emph{Not said in class.} Show that for $f\in C^\infty(M)$ one has $\mathrm d_pf=\partial_a(f\circ x^{-1})(x(p))\,\mathrm dx^a|_p$, by evaluating both sides on $\partial/\partial x^b|_p$ (\cref{thm:tb-dualbasis,def:tv-ddx}).
\end{exercise}

% ==================================================================
\section{The push-forward}\label{sec:tb-pushforward}

The next short section is a very important one, with two important concepts: the \textbf{push-forward} and its cousin, the \textbf{pull-back} \ts{10}{1}{30:13} \ts{10}{1}{30:34} \ts{10}{1}{31:01}.
\begin{definition}[New concepts: push-forward]\label{def:tb-push}
Let $\phi\colon M\to N$ be a smooth map between smooth manifolds (each with its own topology and atlas, a priori unrelated; smoothness of $\phi$ is defined by pushing $\phi$ down to the chart representations, \cref{def:df-map}) \ts{10}{1}{31:01} \ts{10}{1}{31:24} \ts{10}{1}{31:45}.
The \textbf{push-forward} of $\phi$ at $p\in M$ is the linear map
\[
  \phi_{*,p}\colon T_pM\longrightarrow T_{\phi(p)}N,\qquad\bigl(\phi_{*,p}(X)\bigr)(f):=X(f\circ\phi)\quad\text{for }f\in C^\infty(N).
\]
\end{definition}
(In class the symbol was first written with the star up, $\phi^*$, and corrected to the standard convention $\phi_*$ with the star down; ``it's nothing to do with the dual star''. At one point also the target is said to be a ``cotangent space'', a slip: it is $T_{\phi(p)}N$.)
Why is the definition forced? The result must be an element of the tangent space of the target, hence it must act on a smooth function $f$ on $N$; but $f\circ\phi$ is a smooth function \emph{on $M$}, and $X$, a tangent vector to $M$ at $p$, can act on it. The only point of $N$ that is distinguished, given $p\in M$ and a map $\phi$, is $\phi(p)$; the only
linear map from $T_pM$ to $T_{\phi(p)}N$ one can build from the given data is this one. ``Don't memorise stuff you don't need to memorise; every time I construct this when I need it'' \ts{10}{1}{31:45} \ts{10}{1}{33:11} \ts{10}{1}{34:34} \ts{10}{1}{35:36} \ts{10}{1}{36:59} \ts{10}{1}{38:24}.
(Linearity of $\phi_{*,p}$ in $X$ follows from that of $X$; one may also check that $\phi_{*,p}(X)$ is linear and satisfies the Leibniz rule, i.e.\ is a derivation, \cref{def:tv-derivation}; not carried out in class.)

\begin{remark}[Names: derivative, push-forward]\label{rem:tb-names}
(a) As said, it is the only linear map $T_pM\to T_{\phi(p)}N$ that can be constructed from the given data. (b) A small $p$ may be written for each point (``a different such map at every point''): $\phi_{*,p}$. The push-forward is labelled by the star \emph{down} by convention. (c) $\phi_{*,p}$ is often called the \textbf{derivative} of $\phi$ at $p$: when somebody
talks about the derivative of a smooth map between smooth manifolds, this induced map from tangent space to tangent space is meant. (d) \emph{Push-forward} is a very good name: it pushes vectors from the tangent space in the domain to the tangent space in the target, in the direction in which the map goes. If $\phi$ is a diffeomorphism one could also go
back with the inverse, but maps are often not invertible and may nevertheless be smooth in one direction: even then every vector in the domain can be pushed to a vector in the target \ts{10}{1}{38:24} \ts{10}{1}{39:10} \ts{10}{1}{39:31} \ts{10}{1}{40:12} \ts{10}{1}{40:32} \ts{10}{1}{41:13}.
\end{remark}
\begin{proposition}[Push-forward of the tangent vector of a curve]\label{prop:tb-push-curve}
For a smooth map $\phi\colon M\to N$ and a smooth curve $\gamma$ through $p$ (with $\gamma(0)=p$)
\[
  \phi_{*,p}\bigl(X_{\gamma,p}\bigr)=X_{\phi\circ\gamma,\phi(p)}:
\]
the tangent vector of $\gamma$ at $p$ is pushed forward to the tangent vector of the curve $\phi\circ\gamma$ at $\phi(p)$ \ts{10}{1}{41:33} \ts{10}{1}{41:55} \ts{10}{1}{42:17} \ts{10}{1}{42:37}.
\end{proposition}
\emph{Illustration.} Two manifolds $M,N$; in $M$ a curve $\gamma$ through $p$ and its tangent vector (``I do something terrible: I draw a little arrow, but you know it is the directional derivative operator; in charts it corresponds to a $\dim M$-tuple of numbers, the components'', allowed for illustration only). The map takes every point to a point, in a smooth manner,
hence also the points of the curve: the image curve $\phi\circ\gamma$ in $N$, passing through $\phi(p)$; one can construct its tangent vector $Y$ directly, without pushing forward. The claim is that $Y=\phi_{*,p}(X)$ \ts{10}{1}{43:03} \ts{10}{1}{43:43} \ts{10}{1}{44:04} \ts{10}{1}{44:48} \ts{10}{1}{45:09} \ts{10}{1}{45:51}.
\begin{proof}
``It is good practice for the formalism.'' Let $f\in C^\infty(N)$ be arbitrary (a function on the target, since the result is a vector there). By the definition of the push-forward and then of the tangent vector of a curve (with $\gamma(0)=p$; otherwise put the parameter value where $\gamma$ passes through $p$),
\[\begin{aligned}
  \bigl(\phi_{*,p}X_{\gamma,p}\bigr)(f)&=X_{\gamma,p}(f\circ\phi)\\
  &=\bigl((f\circ\phi)\circ\gamma\bigr)'(0)\\
  &=\bigl(f\circ(\phi\circ\gamma)\bigr)'(0)\\
  &=X_{\phi\circ\gamma,\phi(p)}(f),
\end{aligned}\]
using the associativity of the composition and the definition of the tangent vector backwards: $\phi\circ\gamma$ is a curve in $N$ through $\phi(\gamma(0))=\phi(p)$ and $f$ a function on $N$. The two operators agree on an arbitrary function, hence they are equal \ts{10}{1}{45:51} \ts{10}{1}{46:33} \ts{10}{1}{47:35} \ts{10}{1}{48:17} \ts{10}{1}{49:20} \ts{10}{1}{50:01}.
\end{proof}
This is the geometric notion and intuition of the push-forward, ``or, as it's called, the derivative of the function at that particular chosen point'', and it can be done at every point \ts{10}{1}{50:01} \ts{10}{1}{50:22}.
\begin{exercise}[Chain rule for the push-forward]\label{ex:tb-chain}
\emph{Not said in class.} Show that for smooth $\phi\colon M\to N$, $\psi\colon N\to P$ one has $(\psi\circ\phi)_{*,p}=\psi_{*,\phi(p)}\circ\phi_{*,p}$ and $(\mathrm{id}_M)_{*,p}=\mathrm{id}_{T_pM}$ (hence for a diffeomorphism $\phi_{*,p}$ is invertible). Use \cref{def:tb-push}.
\end{exercise}

% ==================================================================
\section{The pull-back}\label{sec:tb-pullback}

\begin{definition}[New concepts: pull-back]\label{def:tb-pull}
Let $\phi\colon M\to N$ be again a smooth map. The \textbf{pull-back} of $\phi$ at the point $\phi(p)\in N$ is the linear map
\[
  \phi^{*}_{p}\colon T^*_{\phi(p)}N\longrightarrow T^*_pM,\qquad\bigl(\phi^*_p(\omega)\bigr)(X):=\omega\bigl(\phi_{*,p}(X)\bigr)\quad\text{for }X\in T_pM,
\]
this time with the star up (\cref{def:tb-push}) \ts{10}{1}{50:47} \ts{10}{1}{51:11} \ts{10}{1}{51:32} \ts{10}{1}{51:53} \ts{10}{1}{52:13} \ts{10}{1}{52:35}.
\end{definition}
It looks like the push-forward, only the cotangent spaces instead of the tangent spaces; the big issue is that the arrow goes in the \emph{opposite direction}, hence the name \emph{pull-back}: if $\phi$ is a diffeomorphism, one can go back and forth and write the pull-back of the inverse, which goes in the same direction; but this is not always the case: $\phi$
need not be surjective, for instance if $M$ has higher dimension than $N$ it cannot be injective and has no inverse. (In class the arrow was written ``for dramatic effect'' in the unusual direction; the point is as stated.) The map was a homework problem but ``it's too important not to do it right now'', and it was constructed together with the students
\ts{10}{1}{52:55} \ts{10}{1}{53:16} \ts{10}{1}{53:37} \ts{10}{1}{53:57}.
\emph{Construction, with the students.} Take $\omega\in T^*_{\phi(p)}N$, a co-vector on the target; $\phi^*_p(\omega)$ should be an element of $T^*_pM$, i.e.\ it is defined by its action on an arbitrary $X\in T_pM$ giving a real number, \emph{linearly} in $X$. The data are only $\phi$, $X$ and $\omega$. A student suggests: use the push-forward, which takes
$X\in T_pM$ to $\phi_{*,p}(X)\in T_{\phi(p)}N$. Then it remains to let $\omega$ act on it (``you said we let this act on $\omega$; no, the other way round: $\omega$ is an element of $T^*_{\phi(p)}N$ so it can act on this vector, it produces a real number''). It is linear because the push-forward is linear and $\omega$ is linear: sums and factors
pull out. Again it is the only way to define it, and one needs only to know from where to where it maps, and that it is linear \ts{10}{1}{53:57} \ts{10}{1}{54:58} \ts{10}{1}{55:40} \ts{10}{1}{56:42} \ts{10}{1}{57:25} \ts{10}{1}{58:27}.
\begin{intuition}[Motto: vectors are pushed forward, co-vectors are pulled back]\label{int:tb-motto}
Memorise: \emph{vectors are pushed forward, forms (co-vectors, and extensions, to come later) are pulled back}. Only under the special circumstance that $\phi$ is a diffeomorphism can one go back and forth, and then one can also pull vectors back, using the inverse map; but if the inverse does not exist, vectors can generically only be pushed forward and co-vectors
only be pulled back. This is very important and will come back once immersions and embeddings are introduced \ts{10}{1}{58:27} \ts{10}{1}{58:51} \ts{10}{1}{59:17} \ts{10}{1}{59:37}.
\end{intuition}
\begin{exercise}[Pull-back of a gradient]\label{ex:tb-pullback-df}
\emph{Not said in class.} Show that $\phi^*_p(\mathrm d_{\phi(p)}f)=\mathrm d_p(f\circ\phi)$ for $f\in C^\infty(N)$ (\cref{def:tb-gradient,def:tb-pull}).
\end{exercise}

% ==================================================================
\section{Immersions and embeddings}\label{sec:tb-immersion}

The push-forward can be used to decide the question: under which circumstances can a smooth manifold ``sit'' (quotation marks) inside some $\R^N$? \ts{10}{1}{1:00:00} \ts{10}{1}{1:00:28} \ts{10}{1}{1:00:49}. There are two notions of sitting, and $\R^N$ is chosen because it is immediately intuitive; one could generalise to one manifold sitting in another: \textbf{immersion} and
\textbf{embedding}; the embedding is the stronger one: every embedding is an immersion, but an immersion has fewer conditions, and sometimes one is happy with an immersion, sometimes one wants an embedding \ts{10}{1}{1:01:09} \ts{10}{1}{1:01:30} \ts{10}{1}{1:01:51}.
\begin{definition}[New concepts: immersion]\label{def:tb-immersion}
A smooth map $\phi\colon M\to\R^N$ is an \textbf{immersion} of $M$ into $\R^N$ if its derivative is injective at every point: the push-forward $\phi_{*,p}\colon T_pM\to T_{\phi(p)}\R^N$ is injective for all $p\in M$. (Injective: two different vectors are not mapped to the same target vector.)
\ts{10}{1}{1:02:13} \ts{10}{1}{1:02:33} \ts{10}{1}{1:02:56} \ts{10}{1}{1:03:18} \ts{10}{1}{1:03:41} \ts{10}{1}{1:04:02}.
\end{definition}
\begin{example}[Two immersions of the circle]\label{ex:tb-circle}
The technical checks are left as homework; intuitively:
\begin{itemize}
  \item \emph{(a)} $\phi\colon S^1\to\R^2$ takes the circle to a distorted circle in the plane. It is smooth (no corners are produced) and its push-forward is injective: the tangent space of the circle is one dimensional and its vectors, pointing in this or that direction, are mapped to vectors that do not coincide.
        (One could also immerse the circle into $\R^3$, as a loop in space, as long as there are no corners.)
  \item \emph{(b)} $\tilde\phi\colon S^1\to\R^2$ takes the circle to a figure-eight, a Lissajous-like curve (a sine in one direction and a sine with double frequency in the other). It is smooth, and \emph{not injective} as a map of sets: two points of the circle are mapped to the crossing point of the figure, one on each branch. But injectivity of
        the map was not required: the derivative $\tilde\phi_{*,p}$ is injective at every point, since the vectors at the two preimages of the crossing are mapped to different vectors, and no vector is mapped to the same as another. So $\tilde\phi$ is an immersion; ``an immersion is not such a strong condition''
\end{itemize}
\ts{10}{1}{1:04:22} \ts{10}{1}{1:05:25} \ts{10}{1}{1:06:08} \ts{10}{1}{1:07:09} \ts{10}{1}{1:07:51} \ts{10}{1}{1:08:54}.
\end{example}
\begin{definition}[New concepts: embedding]\label{def:tb-embedding}
A smooth map $\phi\colon M\to N$ is an \textbf{embedding} if (a) $\phi$ is an immersion and (b) $M$ is homeomorphic to its image $\phi(M)$, which is required to be a submanifold of the target $N$. (The lecture says: if one maps the entire domain, the image should be topologically isomorphic to the manifold; ``if it is topologically isomorphic and smooth,
it immediately implies that it is also diffeomorphic'', but one only needs to check that the image is again a manifold, a submanifold of the target.) It is then not difficult to see that every embedding is an immersion, by condition (a) \ts{10}{1}{1:08:54} \ts{10}{1}{1:09:19} \ts{10}{1}{1:09:40} \ts{10}{1}{1:10:01} \ts{10}{1}{1:10:22}.
\end{definition}
\begin{coursenote}[Wording of the embedding condition]
The usual definition asks that $\phi$ be a homeomorphism onto its image with the subspace topology, equivalently a diffeomorphism onto an embedded submanifold; the lecture's phrasing (``homeomorphic to the target'', corrected to the image) is loose, and the claim that smooth plus one-to-one implies diffeomorphic is not a general fact
(think of $t\mapsto t^3$, a smooth homeomorphism of $\R$ that is not a diffeomorphism), but for an immersion that is a homeomorphism onto its image it is true. Not said in class.
\end{coursenote}
So example (a) is an embedding: the image, the distorted circle, is again a manifold (a smooth one, since $\phi$ is smooth), ``the idea of an embedding: one can think of the smooth manifold as an immersed submanifold of the bigger space''; whereas example (b) is an immersion that is \emph{not} an embedding: at the crossing point the image is not locally homeomorphic to $\R^1$, it fails to be a
manifold. Sometimes one insists on this being a manifold \ts{10}{1}{1:10:22} \ts{10}{1}{1:10:43} \ts{10}{1}{1:11:04} \ts{10}{1}{1:11:26} \ts{10}{1}{1:11:47}.

% ==================================================================
\section{The Whitney theorems}\label{sec:tb-whitney}

So what is the answer to the original question: under which circumstances can a given manifold sit in $\R^N$ for a suitable $N$? \ts{10}{1}{1:11:47} \ts{10}{1}{1:12:08}
\begin{theorem}[Whitney embedding and immersion theorem]\label{thm:tb-whitney}
Any smooth manifold $M$ can be embedded in $\R^{2\dim M}$ and immersed in $\R^{2\dim M-1}$ \ts{10}{1}{1:12:08} \ts{10}{1}{1:12:39} \ts{10}{1}{1:12:59} \ts{10}{1}{1:13:22} \ts{10}{1}{1:13:43} \ts{10}{1}{1:15:25}.
\end{theorem}
(In class the order was confused: first ``it can be embedded in $\R^{2\dim M}$'', then the immersion ``in $2\dim M-1$'' after the Klein bottle example, once the stronger notion was recalled; the statement above is the one the lecture arrives at, and is the correct one.) The good news: there is no smooth manifold that cannot sit in some $\R^N$. The embedding space must
have in the worst case up to twice the dimension of the manifold; some manifolds can be embedded in $\R^{2\dim M-3}$ or the like, it depends on the manifold, but the worst case, ``I think provided by projective spaces'' [?], needs the double dimension \ts{10}{1}{1:13:43} \ts{10}{1}{1:14:03}.
\begin{example}[Circle, surfaces, Klein bottle]\label{ex:tb-klein}
The circle is one dimensional, so it can be embedded in $\R^2$: no surprise. Any two-dimensional manifold (sphere, torus, \dots) can be embedded in $\R^4$. A two-dimensional manifold that cannot be embedded in $\R^3$, which was on the blackboard before: the \textbf{Klein bottle}; but it can be immersed in $\R^{2\cdot2-1}=\R^3$. The picture: a throat
coming in, passing through the wall of the bottle and joining the other side (``like taking a torus and not gluing it like this but going in somewhere else and joining it from the other side''); in $\R^3$ it intersects itself, like the figure-eight before, which prevents the image from being a manifold, hence from being homeomorphic to the Klein bottle: an immersion, not an embedding.
So Whitney guarantees an embedding in $\R^4$, and the Klein bottle is an example in which one cannot beat that limit; for other manifolds one can obviously beat it: the two-sphere can be embedded in $\R^3$ (in the lecture ``immersed into $\R^4$ for sure, but also into $\R^3$'')
\ts{10}{1}{1:14:03} \ts{10}{1}{1:15:05} \ts{10}{1}{1:16:07} \ts{10}{1}{1:17:11} \ts{10}{1}{1:18:14} \ts{10}{1}{1:19:17}.
\end{example}
\begin{coursenote}[Slips in the statement]
At \ts{10}{1}{1:15:25} the lecture first gives ``embedded in $2\dim M-1$'' for the Klein bottle and then corrects to the immersion; and at \ts{10}{1}{1:22:48}--\ts{10}{1}{1:23:08} ``the strong Whitney theorem says a three-dimensional manifold can be embedded in $\R^5$'' where $2\cdot3-1=5$ is the immersion bound and the embedding bound is $\R^6$. The two-sphere is also in the lecture
described as immersed where embedded is meant. Not said in class.
\end{coursenote}
\begin{secondary}[Weaker, stronger and sharper versions]
The statement above is the \emph{strong} Whitney theorem, which takes some effort to prove; one can prove weaker versions more quickly, for example that any manifold can be embedded in \emph{some} $\R^N$ (``also that is not obvious from the beginning'') \ts{10}{1}{1:19:17} \ts{10}{1}{1:19:37} \ts{10}{1}{1:19:58}.
There are even stronger versions that improve the dimension bound: any smooth manifold $M$ can be \emph{immersed} into $\R^{2\dim M-\alpha(\dim M)}$, where $\alpha(n)$ is the number of ones in the binary expansion of the integer $n\ne0$. Example: $\dim M=3$ is $11$ in binary ($1\cdot2^1+1\cdot2^0$), so $\alpha(3)=2$ and any three-dimensional manifold can be immersed in $\R^{6-2}=\R^4$,
whereas the strong Whitney theorem gives $\R^5$. ``I find it quite amusing.'' In a particular case with even better estimates needed one should go to the literature: there are other bounds under particular additional circumstances which may or may not apply
\ts{10}{1}{1:19:58} \ts{10}{1}{1:20:40} \ts{10}{1}{1:21:22} \ts{10}{1}{1:22:26} \ts{10}{1}{1:23:08} \ts{10}{1}{1:23:53}.
\end{secondary}
\begin{history}[Whitney and the immersion bound]
\emph{Not said in class.} \emph{Checked.} H.~Whitney proved the embedding theorem in $\R^{2n}$ and the immersion theorem in $\R^{2n-1}$ in 1936 (\emph{Differentiable manifolds}, \emph{Ann.\ Math.}\ 37) and the easy embedding in $\R^{2n+1}$ earlier; the refinement to $\R^{2n-\alpha(n)}$ for immersions is the theorem of R.~Cohen (1985), the Whitney--Massey
immersion conjecture.\par
\emph{References.} H.~Whitney, \emph{Ann.\ Math.}\ 37 (1936) 645--680; R.~L.~Cohen, \emph{Ann.\ Math.}\ 122 (1985) 237--328.
\end{history}
The conclusion of the lecture: the Whitney theorem \emph{clarifies} whether the intrinsic approach is more powerful than the extrinsic one: it is not, they are the same. One could now check, although it was not proved, that the tangent spaces constructed intrinsically are vector space isomorphic to the tangent planes (higher dimensional planes, not hyperplanes) one would construct for the embedded manifold. For
this one needs an embedding (an immersion alone would not tell how to attach the tangent plane), but one goes one dimension higher at worst. The two approaches are equivalent; this was argued for a while in the history of differential geometry; the intrinsic approach is more aesthetic, and it also justifies drawing the tangent spaces as one intuitively draws them
\ts{10}{1}{1:23:53} \ts{10}{1}{1:24:14} \ts{10}{1}{1:24:35} \ts{10}{1}{1:24:55} \ts{10}{1}{1:25:16} \ts{10}{1}{1:25:36}.
\begin{recall}[Why intrinsic]
The reasons for the intrinsic construction were given in \cref{int:df-embedding}: no reference to an ambient space.
\end{recall}

% ==================================================================
\section{The tangent bundle}\label{sec:tb-bundle}

Another technically highly important topic: the \textbf{tangent bundle}, associated with which one looks at vector fields \ts{10}{1}{1:25:36} \ts{10}{1}{1:25:56}. So far one looked at a manifold and at the tangent space at one particular point; ``after the Whitney embedding theorem I feel more justified in drawing such childish pictures'' \ts{10}{1}{1:25:56} \ts{10}{1}{1:26:17}.
\begin{definition}[New concepts: tangent bundle, bundle projection]\label{def:tb-bundle}
Let $M$ be a smooth manifold. The \textbf{tangent bundle} of $M$ is the set
\[
  TM:=\bigsqcup_{p\in M}T_pM,
\]
the \emph{disjoint} union of all the tangent spaces: every tangent vector at every point of the manifold, as a set. The \textbf{bundle projection} is the map
\[
  \pi\colon TM\to M,\qquad X\mapsto p\qquad\text{where }X\in T_pM,
\]
sending a tangent vector to its base point \ts{10}{1}{1:26:38} \ts{10}{1}{1:26:59} \ts{10}{1}{1:29:05} \ts{10}{1}{1:29:25} \ts{10}{1}{1:30:07} \ts{10}{1}{1:30:29}.
\end{definition}
In pictures the tangent spaces $T_pM$ and $T_qM$ at different points seem to intersect, but one is no longer fooled: that is an artifact of the embedding; intrinsically no element of $T_qM$ is an element of $T_pM$. The dot on the union sign is just an emphasis. The fact that $\pi$ is a \emph{map} (one base point for each $X$) is a consequence of the disjointness
\ts{10}{1}{1:27:41} \ts{10}{1}{1:28:01} \ts{10}{1}{1:28:22} \ts{10}{1}{1:28:43} \ts{10}{1}{1:30:49} \ts{10}{1}{1:31:10}.
So one has a set $TM$ and a map $TM\to M$; this already looks like a bundle (\cref{def:mf-bundle}), but at this point only a \emph{bundle of sets}, not even a topological space bundle, let alone a topological or smooth manifold bundle \ts{10}{1}{1:31:10} \ts{10}{1}{1:31:31} \ts{10}{1}{1:31:52}.

\begin{remark}[A correction on the example of a non-fibre bundle]\label{rem:tb-correction}
In an earlier lecture the lecturer was asked for an example of a bundle that is not a fibre bundle, i.e.\ where not every fibre is the same as the others, and came up quickly with a manifold with a complex line attached at some points and real lines at the others. As a bundle of \emph{sets}, with the projection $\pi$ on top, it is clearly a bundle (no extra
requirement on $\pi$: one cannot talk about continuity without a topology on both sides) and not a fibre bundle: the fibres $\pi^{-1}(p)$ and $\pi^{-1}(q)$ are not isomorphic as sets, there is no bijection between $\R$ and $\mathbb C$ [?: ``that's tricky, there are proofs that you can, well certainly you can get some fibres here and some fibres there'', the lecturer
realises the cardinalities of $\R$ and $\mathbb C$ are equal]. The emphasis is on the way bundles were defined later, as bundles of \emph{topological manifolds} (\cref{def:mf-fibre-bundle}), where continuity is needed; and the biggest problem is that the total space is not a manifold: if $M$ is one dimensional, at one point the total space would be real two dimensional and at another real
three dimensional. So the example is one of a bundle of sets, but not of manifolds; it is really tricky to construct a non-fibre bundle for manifolds, ``but I think it's possible''. ``That was my sin there''; one should not take the example too seriously, although as a bundle of sets it can be held up \ts{10}{1}{1:32:14} \ts{10}{1}{1:32:56} \ts{10}{1}{1:33:16} \ts{10}{1}{1:33:58} \ts{10}{1}{1:34:19} \ts{10}{1}{1:34:59}.
\end{remark}
\begin{coursenote}[Sets of the same cardinality]
Indeed $\R$ and $\mathbb C$ are in bijection as sets (both have the cardinality of the continuum, \cref{ch:l03}); the example is therefore an example of a bundle whose fibres have different \emph{vector space} dimension, not different cardinality. Not said in class.
\end{coursenote}

\paragraph{From a bundle of sets to a smooth manifold.} How to make this into a bundle of topological manifolds, or of smooth manifolds? The base $M$ is smooth. The simplest would be to take the smooth atlas on $M$ and construct from it a smooth atlas on $TM$: then $TM$ is a manifold. A topology on $TM$: one wants that $\pi$ be continuous; the topology to establish is the
\textbf{initial topology} with respect to $\pi$ (it was on the last problem sheet; recall \cref{ch:l04}); in the picture, that already gives open sets of the form $\pi^{-1}(U)$ \ts{10}{1}{1:34:59} \ts{10}{1}{1:35:20} \ts{10}{1}{1:35:41} \ts{10}{1}{1:36:22} \ts{10}{1}{1:36:43} \ts{10}{1}{1:39:51}.
Idea: construct a smooth atlas on $TM$ from a given smooth atlas $\mathcal A_M$ on $M$. The picture: the manifold with the tangent space attached at every point, which looks funny since the tangent space should be tangent, but one does not need to subscribe to the embedded notion, so one can make a nice picture by putting the tangent spaces up; the whole is $TM$
\ts{10}{1}{1:37:06} \ts{10}{1}{1:37:26} \ts{10}{1}{1:37:47}.
Take a chart $(U,x)\in\mathcal A_M$. The new chart domain is the \emph{pre-image} $\pi^{-1}(U)$ (written $\pi^{-1}$ as usual, ``but that has always this taste of $\pi$ being invertible, which it is not at all''), an open subset of $TM$ in the initial topology \ts{10}{1}{1:37:47} \ts{10}{1}{1:38:07} \ts{10}{1}{1:38:28} \ts{10}{1}{1:39:10} \ts{10}{1}{1:39:30} \ts{10}{1}{1:40:12}.
The new chart map is $\xi\colon\pi^{-1}(U)\to\R^{2\dim M}$ (in the lecture the symbol sounds like ``chi''), a subset of $\R^{2\dim M}$ as image; it is intuitively clear that the dimension doubles, since the tangent space has the same dimension as $M$ and ``locally it is trivial'' \ts{10}{1}{1:40:12} \ts{10}{1}{1:40:32} \ts{10}{1}{1:40:52} \ts{10}{1}{1:41:12} \ts{10}{1}{1:41:33}.
\begin{definition}[New concepts: the induced chart on $TM$]\label{def:tb-chart}
For $X\in\pi^{-1}(U)$ (a point of the tangent bundle, an element of some tangent space) define
\[
  \xi(X):=\Bigl(x^1(\pi(X)),\dots,x^{\dim M}(\pi(X)),\ X^1,\dots,X^{\dim M}\Bigr)\in\R^{2\dim M},
\]
where the $X^a$ are the components in the expansion $X=X^a\,\dfrac{\partial}{\partial x^a}\Big|_{\pi(X)}$ with respect to the basis of the tangent space $T_{\pi(X)}M$ induced by the chart $(U,x)$. The first $\dim M$ coordinates are the coordinates of the base point $\pi(X)\in U$ (project down and use the
original chart: they say in which fibre $X$ lies); the other $\dim M$ say where in the fibre $X$ lies. This is possible because $X$ lies in the tangent space over its own base point, and $\partial/\partial x^a|_{\pi(X)}$ is a basis of it; the $X^a$ are real numbers
\ts{10}{1}{1:41:33} \ts{10}{1}{1:42:35} \ts{10}{1}{1:43:59} \ts{10}{1}{1:45:03} \ts{10}{1}{1:46:25} \ts{10}{1}{1:47:28}.
\end{definition}
This is a chart that assigns to an arbitrary point of the set $\pi^{-1}(U)$ a collection of $2\dim M$ numbers. What remains to be shown: that any two such charts $(\pi^{-1}(U),\xi)$ and $(\pi^{-1}(\tilde U),\tilde\xi)$ are $C^\infty$ compatible; this goes on the problem sheet (``quite simple''). Once this is shown, a smooth structure on $M$ is inherited by $TM$, which
is then a smooth manifold; one can speak of the projection $\pi$ being smooth, ``and then you have a smooth bundle'' \ts{10}{1}{1:47:28} \ts{10}{1}{1:47:50} \ts{10}{1}{1:48:12} \ts{10}{1}{1:48:32}.
\begin{exercise}[The charts of $TM$ are compatible]\label{ex:tb-compat}
Show that the transition map $\tilde\xi\circ\xi^{-1}$ between two induced charts is $C^\infty$: the first $\dim M$ components are $\tilde x\circ x^{-1}$, and the last ones are the Jacobian of $\tilde x\circ x^{-1}$ at the base point applied to $(X^a)$ (\cref{ex:tv-change}), linear in $X^a$. (Problem sheet; the full computation, not given in class, is in
\cref{exam:ch10}.)
\end{exercise}

% ==================================================================
\section*{Key concepts}
\begin{keyconcepts}
\textbf{Retrieval questions.}
\begin{enumerate}
  \item Define $T^*_pM$ and $T^r_s(T_pM)$. Why is $T_pM\cong T^*_pM$ non-canonical, while $V\cong V^{**}$ is canonical?
  \item Define the gradient $\mathrm d_pf$. Why is it a co-vector, and why does ``the gradient is orthogonal to the level sets'' not need an inner product?
  \item Prove that $\mathrm dx^a|_p$ is the dual basis of $\partial/\partial x^b|_p$. What does $\mathrm dx^a$ \emph{not} mean?
  \item Define the push-forward $\phi_{*,p}$ and show that it is forced by the data; prove $\phi_{*,p}X_{\gamma,p}=X_{\phi\circ\gamma,\phi(p)}$.
  \item Define the pull-back. Why does it go in the opposite direction? State the motto.
  \item Define immersion and embedding; give an immersion of the circle that is not an embedding.
  \item State the Whitney theorems, with the Klein bottle as an example and the improved bound $2\dim M-\alpha(\dim M)$. What do they say about the intrinsic versus the extrinsic definition of tangent spaces?
  \item Define the tangent bundle, the bundle projection and the chart on $TM$ induced by a chart on $M$. What remains to be proven?
\end{enumerate}
\textbf{Concepts.}
\begin{itemize}
  \item Cotangent space, tensors at a point (\cref{def:tb-cotangent,def:tb-tensors,rem:tb-noncanonical}); gradient (\cref{def:tb-gradient,rem:tb-covector}); chart-induced dual basis (\cref{def:tb-dx,thm:tb-dualbasis,int:tb-dx}).
  \item Push-forward and pull-back (\cref{def:tb-push,prop:tb-push-curve,def:tb-pull,int:tb-motto}).
  \item Immersion, embedding, Whitney (\cref{def:tb-immersion,def:tb-embedding,ex:tb-circle,thm:tb-whitney,ex:tb-klein}).
  \item Tangent bundle and its smooth structure (\cref{def:tb-bundle,rem:tb-correction,def:tb-chart,ex:tb-compat}).
\end{itemize}
\end{keyconcepts}

\flushside
\end{paracol}
